\section{Conclusiones}

Se diseñó e implementó una base de datos relacional para procesos odontológicos, cubriendo registro de pacientes, citas, odontogramas, tratamientos, pagos, compras e inventario.

El modelo conceptual Chen permitió representar el dominio sin depender de un motor físico. La transformación relacional permitió derivar tablas, claves y relaciones. El modelo lógico Crow's Foot facilitó la visualización de la estructura relacional, mientras que Oracle Database 19c permitió implementar el diseño físico con restricciones, vista, triggers y procedimiento almacenado.

Las validaciones realizadas demostraron que los objetos fueron creados correctamente y que las reglas implementadas funcionan. En particular, se verificó la actualización de inventario al registrar compras y la actualización de saldo al registrar pagos mediante el procedimiento \texttt{SP\_REGISTRAR\_PAGO}.

Finalmente, la comparación con el enfoque MySQL/Laravel descrito en la tesis base permitió evidenciar que Oracle 19c ofrece una implementación más centrada en reglas del motor, integridad relacional explícita y validación técnica mediante el catálogo de base de datos.


La ejecución del repositorio Laravel del proyecto permitió comprobar una aplicación funcional del mismo dominio odontológico usando Oracle, migraciones, seeders y autenticación. Esta verificación no sustituye la base diseñada en el examen; más bien, refuerza la trazabilidad entre el sistema web del proyecto LALYSDENT y el diseño independiente implementado en el schema \texttt{FARMACIAS\_BD3}.

El uso de schemas separados permitió mantener una distinción clara entre \texttt{FARMACIAS\_BD3}, como producto principal del examen, y \texttt{LALYSDENT\_APP}, como entorno de ejecución de la aplicación Laravel desarrollada por el equipo.
